\chapter{System Architecture at a Glance}
\label{app:architecture}

This appendix is a map. It gives a high-level tour of how Mathilda is put
together---the handful of core pieces every computation flows through---and points
at where in the source tree each piece lives, so that a curious reader can open the
right file. It is deliberately compact; the authoritative engineering reference is
\texttt{SPEC.md} at the root of the source tree, and Chapter~\ref{ch:internals}
develops these ideas in depth.

\section{The pipeline}
\label{app:pipeline}

At the largest scale Mathilda is a loop that turns a string into an expression,
rewrites that expression until it stops changing, and turns the result back into a
string. Five components do the work, and the evaluator consults the middle three on
every step.

\begin{figure}[h]
\begin{framed}
\small
\begin{verbatim}
    "Integrate[1/(1+x^2), x]"          input string
             |
             v
      [ Parser ]      parse.c      Pratt / operator-precedence
             |
             v   Expr* tree  ->  Integrate[Times[...], x]
             |
      [ Evaluator ]   eval.c       fixed-point loop
        |  ^  ^  ^
        |  |  |  +--- Built-in C library    (plus.c, calculus/, ...)
        |  |  +------ Pattern matcher (match.c) + Rule engine (replace.c)
        |  +--------- Attribute system (attr.c)
        |  +--------- Symbol table (symtab.c)  OwnValues / DownValues
             |
             v
      [ Printer ]     print.c      precedence-aware formatting
             |
             v
        "ArcTan[x]"                  output string
\end{verbatim}
\end{framed}
\caption{The evaluation pipeline. The parser and printer bracket a fixed-point
evaluator that, on every step, consults the symbol table, the attribute system,
the pattern matcher, and the built-in library.}
\label{fig:pipeline}
\end{figure}

The remaining sections walk these pieces in turn.

\section{Everything is an expression}
\label{app:expr}

The single most important design decision in Mathilda is that \emph{everything is an
expression}. A number, a symbol, a list, a matrix, a rule, even a pattern---all are
the same kind of tree, an \texttt{Expr} (defined in \texttt{src/expr.\{c,h\}}). An
\texttt{Expr} is a tagged union: an integer (machine \texttt{int64} or a GMP
bignum), a real, a symbol, a string, or a \emph{function} node carrying a head and a
list of argument expressions. The compound form \texttt{f[x, y]} is a function node
with head \texttt{f}; a list \texttt{\{a, b\}} is really \texttt{List[a, b]}
internally.

You can see the tree behind any surface syntax with \B{FullForm}:

\mtranscript{appendix-architecture/expr-model}

Read those lines as X-rays. The sum \texttt{a + b c} is
\texttt{Plus[a, Times[b, c]]}; a fraction is a \texttt{Rational}; a division is a
\texttt{Times} against a \texttt{Power} with exponent $-1$; and even a pattern like
\texttt{x\_} is an ordinary expression, \texttt{Pattern[x, Blank[]]}. Because atoms
carry a head too---\B{Head} of an integer is \texttt{Integer}, of a list is
\texttt{List}---one uniform set of tools (\B{Part}, \B{Map}, \B{Replace}, and the
rest) operates on \emph{everything}. This uniformity is why the evaluator and the
pattern matcher can be small: they only ever manipulate one data structure.

Integers illustrate a second theme: representations that adapt. A value begins as a
fast machine word and is promoted to a GMP bignum only when an operation would
overflow, demoting back when a result fits again---so there is only ever ``an
integer,'' never a separate big-integer mode.

\section{How a call is evaluated}
\label{app:eval}

Evaluation is a fixed-point process (\texttt{src/eval.c}): \texttt{evaluate} applies
one \texttt{evaluate\_step} after another until the expression stops changing. For a
function call \texttt{f[a1, \dots, aN]} a single step proceeds, in order:

\begin{enumerate}
  \item evaluate the head \texttt{f};
  \item read \texttt{f}'s \textbf{attributes};
  \item evaluate the arguments---except those held back by a \texttt{Hold} attribute;
  \item if \texttt{f} is \textbf{Listable}, thread over any list arguments;
  \item if \texttt{f} is \textbf{Flat}, flatten nested \texttt{f[\dots]} calls;
  \item if \texttt{f} is \textbf{Orderless}, sort the arguments canonically;
  \item call \texttt{f}'s built-in C function, if it has one---this returns a new
        \texttt{Expr*} on success, or \texttt{NULL} to mean ``I can't evaluate this,
        leave it alone'';
  \item otherwise try \texttt{f}'s \textbf{DownValues}---the user- or library-defined
        rewrite rules---and take the first that matches.
\end{enumerate}

The attribute bits (\texttt{src/attr.c}) are the small vocabulary that steers all of
this. They are what let the evaluator stay generic while each head opts into exactly
the behaviour it needs\calloutnote{Under the hood}{\B{Plus} is \texttt{Flat}
(associative), \texttt{Orderless} (commutative), \texttt{Listable} (threads over lists),
\texttt{OneIdentity}, and a \texttt{NumericFunction}; \B{Sin} is merely \texttt{Listable}
and numeric; \B{Table} is \texttt{HoldAll}, so its iterator specification is \emph{not}
evaluated before \texttt{Table} runs. Nothing about the evaluator's main loop changes
between these heads---only their attribute bits do.}:

\mtranscript{appendix-architecture/attributes}

Because a step just rewrites and the loop runs to a fixed point, rule chains compose
on their own, with no explicit plumbing. Define \texttt{f} in terms of \texttt{g} and
\texttt{g} in terms of squaring, and the answer falls out over successive passes:

\mtranscript{appendix-architecture/fixed-point}

\section{Matching, rules, and the built-in library}
\label{app:match}

Two subsystems implement the last two steps above. The \textbf{pattern matcher}
(\texttt{src/match.c}) unifies an expression against a pattern with backtracking,
handling blanks (\texttt{\_}), typed blanks (\texttt{\_h}), sequence patterns
(\texttt{\_\_}, \texttt{\_\_\_}), tests, and conditions, binding pattern variables in
a match environment. The \textbf{rule engine} (\texttt{src/replace.c}) applies
rules---\texttt{Rule} (\texttt{->}, immediate) and \texttt{RuleDelayed}
(\texttt{:>}, delayed)---through \B{Replace}, \B{ReplaceAll} (\texttt{/.}), and
\B{ReplaceRepeated} (\texttt{//.}).

The \textbf{built-in C library} supplies the primitives. Each builtin obeys a single
ownership contract: it takes ownership of its argument expression and returns either
a freshly built \texttt{Expr*} (success) or \texttt{NULL} (``not my case,'' leaving
the expression for a rule to handle). That \texttt{NULL} convention is what lets
symbolic arguments flow untouched through heads that only know how to handle numbers.

\section{The source tree}
\label{app:tree}

The C sources live under \texttt{src/}. The core is a small number of files; the bulk
is the mathematical library, organised so that---as a rule of thumb---each builtin
gets its own file, grouped into a subsystem directory.

\begin{figure}[h]
\begin{framed}
\small
\begin{verbatim}
src/
  expr.{c,h}     expression representation (the tagged union)
  parse.{c,h}    Pratt parser            print.{c,h}   printer
  eval.{c,h}     fixed-point evaluator    symtab.{c,h}  symbol table
  attr.{c,h}     attribute bitflags       match.{c,h}   pattern matcher
  replace.{c,h}  rule engine              core.{c,h}    builtin + init hub
  plus.c times.c power.c arithmetic.c     arithmetic heads
  ndarray.* ndkernels.* pack.*            packed / dense-array substrate
  compile/       Compile[] bytecode compiler and auto-compilation
  calculus/      D, Series, Limit, Integrate, Risch-Norman
  poly/  rat/  simp/  linalg/  solve/     algebra & solving
  numerical_calculus/  numerical_roots/   N-prefixed numerics
  sum/  product/  special_functions/      series, products, transcendentals
  numbertheory/  stats/  graph/  ml/      further subsystems
  strings/  list/  graphics/              text, lists, 2D/3D graphics
  internal/      init.m and .m rule files (derivatives, integral tables)
tests/           CMake-built unit suite
makefile         primary build system
\end{verbatim}
\end{framed}
\caption{A condensed view of the source tree. The full layout, with a line for every
directory, is in \texttt{SPEC.md}~\S2.}
\label{fig:tree}
\end{figure}

\section{Startup and module initialization}
\label{app:init}

When Mathilda starts, \texttt{main()} in \texttt{src/repl.c} calls
\texttt{symtab\_init()} to create the symbol table, then \texttt{core\_init()} in
\texttt{src/core.c}. \texttt{core\_init()} registers the core builtins and then calls
each subsystem's \texttt{*\_init()} function in a fixed order---several dozen of them,
spanning arithmetic, special functions, the packed-array substrate, the numeric
subsystems, and the higher-level algebra and graphics services. \texttt{core.c} is
the source of truth for that order.

After the C side is initialised, \texttt{main()} loads
\texttt{src/internal/init.m}, which in turn \texttt{Get[]}s the remaining bootstrap
\texttt{.m} files. This is the seam of Mathilda's \textbf{two-tier design}:
performance-sensitive primitives are written in C, while higher-level mathematical
identities---the derivative rules in \texttt{deriv.m}, the integral tables---are
expressed as ordinary rewrite rules in Mathilda's own language. The same fixed-point
evaluator that runs a user's rules runs these, so extending the system's mathematics
often means adding a rule, not writing C.

\bigskip
\noindent
For the authoritative detail behind every claim here, see \texttt{SPEC.md}; for
recipes on adding builtins, modules, patterns, and operators, see
\texttt{docs/extending.md} and Chapter~\ref{ch:contributing}.
